\SourceLeaf{406}%
\par\medskip\Needspace{6\baselineskip}
\pdfbookmark[2]{VII.}{carnot-zusatz-vii}
\centerline{\textbf{VII.}}\nobreak\smallskip
\begin{center}\textit{Derivation of the basic formulas of spherical trigonometry.}\end{center}
\hspace*{\parindent}If \(a,b,c\) are the sides of a spherical triangle and \(A,B,C\) the opposite angles, then
\[
\cos a=\cos b.\cos c+\sin b.\sin c.\cos A
\]
as \textsc{Lagrange} elegantly proved when both \(b\) and \(c\) are less than \(90^0\). The proof can readily be extended to all the other cases.

The following cases are possible.
\[
\text{I.}\quad .\ .\ .\ .\ .\ .\qquad b<90^0,\quad c<90^0
\]
\par
Here the proof applies directly.
\[
\text{II.}\quad .\ .\ .\ .\ .\ .\qquad b>90^0,\quad c>90^0
\]
\par
Extend the sides \(b\) and \(c\) beyond the points \(B\) and \(C\) until they meet at \(A'\), and find the value of \(a\) by considering the triangle \(A'BC\).
\[
\text{III.}\quad .\ .\ .\ .\ .\ .\qquad b>90^0,\quad c<90^0
\]
\par
Extend the sides \(b\) and \(a\) beyond the points \(A\) and \(B\) until they meet at \(C'\). In the triangle \(C'AB\), case I then gives
\[
\cos(180^0-a)=\cos c.\cos(180^0-b)+\sin c.\sin(180^0-b).\cos(180^0-A)
\]
which is identical to the basic formula.

The case \(b<90^0\) and \(c>90^0\) is essentially the same.
\par\Needspace{5\baselineskip}
\[
\text{IV.}\quad .\ .\ .\ .\ .\ .\qquad b=90^0,\quad A=90^0
\]
\par
Here \(C\) is the pole of \(AB\), so necessarily \(a=90^0\) as well. The formula
\[
\cos a=\cos b.\cos c+\sin b.\sin c.\cos A
\]
is therefore evident. The case \(c=90^0,\ A=90^0\) is essentially the same.
